%AttackName: M-ScoreBased Maximized-1

%Input:
Let $x$ be the original input image and $y$ be the true label associated with it. The objective is to create an adversarial example $x^*$ that misleads the model.
- $J(f_\theta(x), y)$ is a score function, often based on the logits output by the model,
- $f_\theta(x)$ is the model's prediction for input $x$,
- $\epsilon$ controls the perturbation magnitude.

%Output:
The output of the score-based attack is an adversarial example $x^*$ generated by modifying the original input $x$ based on the scores (logits) provided by the model.

%Formula:
The score-based attack updates the input by maximizing the difference between the target class score and the scores of other classes. This can be expressed as:
$x^* = x + \epsilon \cdot \text{sign}(\nabla_x J(f_\theta(x), y))$

In some formulations, the update can be defined as:
$\text{maximize} \quad \text{score}(f_\theta(x^*), c) - \text{score}(f_\theta(x^*), i)$
for all classes $i \neq c$, where $c$ is the target class.

%Explanation:
The Score-Based Maximized adversarial attack leverages the model's score function, typically derived from the logits output, to craft adversarial examples. By focusing on maximizing the score of a target class while minimizing the scores of other classes, the attack effectively manipulates the input to mislead the model's predictions. The perturbation is constrained by a parameter $\epsilon$, ensuring that the changes to the original input remain within acceptable bounds. This method highlights the vulnerabilities of models that rely heavily on score-based outputs and underscores the importance of robust defenses against such targeted attacks.